\subsubsection{Hamiltonians and Molecular Energy}\label{sec:hamiltonians-and-molecular-energy}

Inside each molecule, there are nuclei and electrons that interact with each other through electromagnetic forces. The question is: \emph{\textbf{how much energy does this molecule have?}} But more importantly: \textbf{\emph{which arrangement of the electrons corresponds to the minimum energy?}}

\highspace
\textcolor{Green3}{\faIcon{question-circle} \textbf{Why can't we compute it directly?}} In principle, we could compute the energy of a molecule directly from the laws of quantum mechanics. However, for realistic molecules this becomes computationally intractable because every particle contributes to the total energy and all particles interact with one another.

\highspace
More precisely:
\begin{itemize}
    \item Each electron has \textbf{kinetic energy};
    \item Each electron interacts with \textbf{every nucleus};
    \item Each electron interacts with \textbf{every other electron};
    \item Nuclei also interact with \textbf{each other}.
\end{itemize}
So the total energy is \textbf{not} simply the \textbf{sum of independent contributions}. It is the result of a \hl{huge number of coupled quantum interactions.}

\highspace
\textcolor{Green3}{\faIcon{check-circle} \textbf{Why introduce the Hamiltonian?}} The Hamiltonian is \textbf{not} introduced because writing many energy terms is inconvenient. It is introduced because \hl{quantum mechanics is formulated in terms of operators acting on quantum states.} In general, the Hamiltonian is the operator that \textbf{contains all the information about the energy of the system} (\autopageref{sec:hamiltonian-recap}, where a system could be a molecule, an atom, a spin chain, a quantum computer, etc). Here, we use the Hamiltonian to describes the \textbf{physical energy of a molecule}; it contains \textbf{all} contributions to the molecular energy.

\highspace
\textcolor{Green3}{\faIcon{question-circle} \textbf{What does the Hamiltonian contain?}} In the context of VQE, the Hamiltonian contains \textbf{all the energy contributions of the molecule}. More precisely, it represents: the \textbf{kinetic energy} of all the particles, and the \textbf{potential energy} of all the particles. Together, these describe the molecule's total energy. The Hamiltonian is therefore a compact mathematical representation of the molecule's entire energy landscape.

\highspace
\begin{flushleft}
    \textcolor{Green3}{\faIcon{tools} \textbf{The Time-Independent Schrödinger Equation}}
\end{flushleft}
Now that we have introduced the operator (Hamiltonian) that contains all the information about the molecule's energy, the next question is: \emph{\textbf{how do we obtain the energies encoded inside the Hamiltonian?}} The answer is given by the \definition{Time-Independent Schrödinger Equation}:
\begin{equation}\label{eq:time-independent-schrodinger-equation}
    H \ket{\psi} = E \ket{\psi}
\end{equation}
The Schrödinger equation is the mathematical tool that \textbf{extracts the energies from the Hamiltonian}. The equation is simply an \textbf{eigenvalue equation}, so it tells us: \emph{find the special states $\ket{\psi}$ for which applying the Hamiltonian simply multiplies the state by a number $E$.} The number $E$ is the \textbf{energy} of the state $\ket{\psi}$.
\begin{itemize}
    \item $H$ is the \textbf{Hamiltonian operator}, which describes the total energy of the molecule;
    \item $\ket{\psi}$ is the \textbf{quantum state of the molecule}, a possible configuration of the electrons. In simple terms, ``\emph{how are the electrons arranged?}'';
    \item $E$ is the \textbf{energy associated} with that particular quantum state $\ket{\psi}$.
\end{itemize}
Solving the Schrödinger equation is equivalent to select the \textbf{special states} that satisfy the Hamiltonian. These states are the \textbf{eigenstates} of the Hamiltonian. Each one has a \textbf{well-defined energy} $E$, which is the corresponding \textbf{eigenvalue}. It is similar to the algebraic case:
\begin{equation*}
    A \vec{v} = \lambda \vec{v}
\end{equation*}
Where $A$ is a matrix, $\vec{v}$ is an eigenvector, and $\lambda$ is the corresponding eigenvalue. In the quantum case, the Hamiltonian is an operator, the quantum state is an eigenvector, and the energy is the corresponding eigenvalue.

\highspace
\textcolor{Green3}{\faIcon[regular]{lightbulb} \textbf{The power of the time-independent formulation.}} With the time-independent Schrödinger equation we are \textbf{not} studying how the state evolves over time. Instead, we are looking for the special states that have a \textbf{well-defined}, \textbf{constant energy}. The solution gives us a state where the energy is fixed and \textbf{does not change into another energy level}. This is why these states are also known as \textbf{stationary states}.

\highspace
In other words, the time-independent Schrödinger equation finds the \textbf{stationary (stable) quantum states} of the Hamiltonian. These states have a \textbf{well-defined energy} that remains \emph{constant over time}. This is why they are considered stable and are the states of interest in VQE.

\highspace
\textcolor{Green3}{\faIcon{question-circle} \textbf{Which stationary state are we looking for?}}
Solving the time-independent Schrödinger equation does not produce a single solution. Instead, it produces a set of \textbf{eigenstates}:
\begin{equation*}
    \ket{\psi_0},\;
    \ket{\psi_1},\;
    \ket{\psi_2},\;\ldots
\end{equation*}
Each associated with an energy eigenvalue
\begin{equation*}
    E_0,\;
    E_1,\;
    E_2,\;\ldots
\end{equation*}
Where, by convention:
\begin{equation*}
    E_0 \le E_1 \le E_2 \le \cdots
\end{equation*}
The state with the \textbf{lowest energy} $\ket{\psi_0}$ is called the \textbf{ground state}. It represents the most stable configuration of the molecule because physical systems naturally tend to occupy the minimum-energy state. All the remaining eigenstates:
\begin{equation*}
    \ket{\psi_1},\;
    \ket{\psi_2},\;\ldots
\end{equation*}
Are called \textbf{excited states} (\autopageref{eq:ground-state-energy}), since they possess higher energies and can only be occupied if additional energy is supplied to the system.

\highspace
\textcolor{Green3}{\faIcon{question-circle} \textbf{Why does VQE only care about these states?}} Because VQE wants to find the \textbf{ground state}, which is \textbf{one of the stationary states}, specifically the one with the \textbf{lowest energy}.

\highspace
\begin{deepeningbox}[: How can we represent the molecular Hamiltonian on a quantum computer?]
    A quantum computer does \textbf{not} understand arbitrary matrices. It understands quantum gates built from Pauli operators. Therefore, before running VQE, the \hl{molecular Hamiltonian is rewritten as a sum of tensor products of Pauli matrices.}

    \highspace
    \textcolor{Green3}{\faIcon{question-circle} \textbf{Why tensor products?}} Because a molecule involves \textbf{multiple qubits}. The tensor product (\autopageref{subsubsec:tensor-product}) is the mathematical operation that combines multiple qubits into a single quantum state. Each qubit can be represented by a Pauli operator, and the tensor product allows us to describe the interactions between these qubits in the Hamiltonian.

    \highspace
    \textcolor{Green3}{\faIcon{question-circle} \textbf{Why a sum?}} Each tensor product represents \textbf{one contribution} to the Hamiltonian. The total Hamiltonian is obtained by adding them together.

    \highspace
    In general, although the Hamiltonian represents the physical energy of the molecule, before it can be processed by a quantum computer it is expressed as a \textbf{sum of tensor products of Pauli matrices}, since Pauli operators form the natural building blocks of quantum circuits.
\end{deepeningbox}

\highspace
In summary, in VQE, the energy of a molecule is represented by the Hamiltonian $H$, a quantum operator that includes the kinetic and potential energies of all the particles in the system. The physical states of the molecule are obtained by solving the time-independent Schrödinger equation $H\ket{\psi} = E\ket{\psi}$, where $\ket{\psi}$ is a quantum state and $E$ is its associated energy. Since this equation has the form of an eigenvalue problem, the solutions correspond to the eigenstates of the Hamiltonian and their associated energy levels.

\highspace
Therefore, the goal of VQE is not to compute all the eigenvalues of the Hamiltonian. Instead, it focuses on finding the \textbf{minimum eigenvalue} and its corresponding eigenstate:
\begin{equation*}
    E_{\min}=E_0,
    \qquad
    \ket{\psi_{\min}}=\ket{\psi_0}
\end{equation*}
Once the ground state has been identified, many important physical and chemical properties of the molecule can be determined.